\documentclass[11pt, a4paper]{report}

% ========== Common Packages ==========
\usepackage{fontspec}
\usepackage{kotex}
\usepackage{amsmath, amssymb, amsthm}
\usepackage{mathtools}
\usepackage{bm}
\usepackage{graphicx}
\usepackage{booktabs}
\usepackage{array}
\usepackage{tabularx}
\usepackage{longtable}
\usepackage{hyperref}
\usepackage{cleveref}
\usepackage{geometry}
\usepackage{fancyhdr}
\usepackage{titlesec}
\usepackage{enumitem}
\usepackage{xcolor}
\usepackage{tcolorbox}
\usepackage{algorithm}
\usepackage{algorithmic}
\usepackage{tikz}
\usetikzlibrary{arrows.meta, positioning, calc, decorations.pathreplacing, shapes.geometric, fit, patterns, angles, quotes, trees}
\usepackage{pgfplots}
\pgfplotsset{compat=1.18}
\usepackage{caption}
\usepackage{subcaption}
\usepackage{float}

\geometry{margin=2.5cm, top=3cm, bottom=3cm}

\usepackage{multirow}
\usepackage{siunitx}
% ========== Colors ==========
% Common
\definecolor{defblue}{RGB}{30, 80, 150}
\definecolor{thmgreen}{RGB}{20, 110, 60}
\definecolor{noteorange}{RGB}{200, 100, 0}
\definecolor{lightblue}{RGB}{230, 240, 255}
\definecolor{lightgreen}{RGB}{230, 255, 235}
\definecolor{lightorange}{RGB}{255, 245, 230}
\definecolor{lightgray}{RGB}{245, 245, 245}
% Extended
\definecolor{lightred}{RGB}{255, 235, 235}
\definecolor{darkred}{RGB}{180, 30, 30}
\definecolor{biogreen}{RGB}{10, 130, 80}
\definecolor{cvpurple}{RGB}{100, 40, 150}
\definecolor{injuryred}{RGB}{200, 40, 40}
\definecolor{codegray}{RGB}{240, 240, 245}
\definecolor{pipeblue}{RGB}{25, 100, 180}
\definecolor{constraintred}{RGB}{200, 50, 50}
\definecolor{termblack}{RGB}{30, 30, 30}
\definecolor{goldcolor}{RGB}{180, 140, 20}

% ========== Theorem-like environments ==========
\theoremstyle{definition}
\newtheorem{definition}{Definition}[chapter]
\newtheorem{example}{Example}[chapter]

\theoremstyle{plain}
\newtheorem{theorem}{Theorem}[chapter]
\newtheorem{proposition}{Proposition}[chapter]

\theoremstyle{remark}
\newtheorem{remark}{Remark}[chapter]

% ========== tcolorbox environments ==========
% --- Common (all parts) ---
\newtcolorbox{keyidea}[1][]{
  colback=lightblue,
  colframe=defblue,
  fonttitle=\bfseries,
  title={Key Idea},
  #1
}

\newtcolorbox{importantnote}[1][]{
  colback=lightorange,
  colframe=noteorange,
  fonttitle=\bfseries,
  title={Important},
  #1
}

\newtcolorbox{summary}[1][]{
  colback=lightgreen,
  colframe=thmgreen,
  fonttitle=\bfseries,
  title={Summary},
  #1
}

% --- Warning/Error ---
\newtcolorbox{warningbox}[1][]{
  colback=lightred,
  colframe=darkred,
  fonttitle=\bfseries,
  title={Warning},
  #1
}

% --- Domain: Biomechanics & Clinical ---
\newtcolorbox{clinicalnote}[1][]{
  colback=lightred,
  colframe=darkred,
  fonttitle=\bfseries,
  title={Clinical Note},
  #1
}

\newtcolorbox{biomechbox}[1][]{
  colback=lightgreen,
  colframe=biogreen,
  fonttitle=\bfseries,
  title={Biomechanics},
  #1
}

\newtcolorbox{injurybox}[1][]{
  colback={injuryred!8},
  colframe=injuryred,
  fonttitle=\bfseries,
  title={Injury Risk},
  #1
}

% --- Domain: Computer Vision ---
\newtcolorbox{cvbox}[1][]{
  colback={cvpurple!5},
  colframe=cvpurple,
  fonttitle=\bfseries,
  title={Computer Vision},
  #1
}

% --- Pipeline & Setup ---
\newtcolorbox{setupbox}[1][]{
  colback=lightgray,
  colframe=gray!60,
  fonttitle=\bfseries,
  title={Setup},
  #1
}

\newtcolorbox{pipelinebox}[1][]{
  colback=pipeblue!5,
  colframe=pipeblue,
  fonttitle=\bfseries,
  title={Pipeline Step},
  #1
}

\newtcolorbox{codebox}[1][]{
  colback=codegray,
  colframe=gray!70,
  fonttitle=\bfseries\ttfamily,
  title={Code},
  #1
}

\newtcolorbox{termbox}[1][]{
  colback=termblack!5,
  colframe=termblack!60,
  fonttitle=\bfseries\ttfamily,
  title={Terminal},
  #1
}

% --- Research ---
\newtcolorbox{hypothesis}[1][]{
  colback=goldcolor!8,
  colframe=goldcolor,
  fonttitle=\bfseries,
  title={Hypothesis},
  #1
}

\newtcolorbox{resultbox}[1][]{
  colback=thmgreen!8,
  colframe=thmgreen,
  fonttitle=\bfseries,
  title={Expected Result},
  #1
}

% --- Constraints & Solutions ---
\newtcolorbox{constraintbox}[1][]{
  colback=constraintred!8,
  colframe=constraintred,
  fonttitle=\bfseries,
  title={Constraint},
  #1
}

\newtcolorbox{solutionbox}[1][]{
  colback=thmgreen!8,
  colframe=thmgreen,
  fonttitle=\bfseries,
  title={Solution},
  #1
}

\newtcolorbox{databox}[1][]{
  colback=pipeblue!5,
  colframe=pipeblue,
  fonttitle=\bfseries,
  title={Data Spec},
  #1
}

% --- Progress/Status ---
\newtcolorbox{successbox}[1][]{
  colback=lightgreen,
  colframe=thmgreen,
  fonttitle=\bfseries,
  title={Success},
  #1
}

\newtcolorbox{nextbox}[1][]{
  colback=lightorange,
  colframe=noteorange,
  fonttitle=\bfseries,
  title={Next Step},
  #1
}

% --- Fitting guide ---
\newtcolorbox{failbox}[1][]{
  colback=lightred,
  colframe=darkred,
  fonttitle=\bfseries,
  title={Failure Pattern},
  #1
}

\newtcolorbox{fixbox}[1][]{
  colback=lightgreen,
  colframe=thmgreen,
  fonttitle=\bfseries,
  title={Fix},
  #1
}

\newtcolorbox{lessonbox}[1][]{
  colback=lightorange,
  colframe=noteorange,
  fonttitle=\bfseries,
  title={Lesson Learned},
  #1
}

% --- Specification ---
\newtcolorbox{specbox}[1][]{
  colback=cvpurple!5,
  colframe=cvpurple,
  fonttitle=\bfseries,
  title={Specification},
  #1
}

\newtcolorbox{checklistbox}[1][]{
  colback=lightgreen,
  colframe=biogreen,
  fonttitle=\bfseries,
  title={Checklist},
  #1
}

% ========== Custom math commands ==========
% Vectors (lowercase)
\newcommand{\vx}{\mathbf{x}}
\newcommand{\vd}{\mathbf{d}}
\newcommand{\vo}{\mathbf{o}}
\newcommand{\vc}{\mathbf{c}}
\newcommand{\vr}{\mathbf{r}}
\newcommand{\vp}{\mathbf{p}}
\newcommand{\vv}{\mathbf{v}}
\newcommand{\vn}{\mathbf{n}}
\newcommand{\vu}{\mathbf{u}}
\newcommand{\vt}{\mathbf{t}}
\newcommand{\ve}{\mathbf{e}}
\newcommand{\vq}{\mathbf{q}}
% Vectors (uppercase)
\newcommand{\vF}{\mathbf{F}}
\newcommand{\vM}{\mathbf{M}}
\newcommand{\va}{\mathbf{a}}
\newcommand{\vg}{\mathbf{g}}
\newcommand{\vX}{\mathbf{X}}
% Bold Greek
\newcommand{\vmu}{\bm{\mu}}
\newcommand{\vbeta}{\bm{\beta}}
\newcommand{\vtheta}{\bm{\theta}}
\newcommand{\vpsi}{\bm{\psi}}
% Matrices
\newcommand{\mSigma}{\bm{\Sigma}}
\newcommand{\mR}{\mathbf{R}}
\newcommand{\mS}{\mathbf{S}}
\newcommand{\mW}{\mathbf{W}}
\newcommand{\mJ}{\mathbf{J}}
\newcommand{\mG}{\mathbf{G}}
\newcommand{\mT}{\mathbf{T}}
\newcommand{\mI}{\mathbf{I}}
\newcommand{\mB}{\mathbf{B}}
\newcommand{\mK}{\mathbf{K}}
\newcommand{\mP}{\mathbf{P}}
\newcommand{\mF}{\mathbf{F}}
\newcommand{\mE}{\mathbf{E}}
\newcommand{\mA}{\mathbf{A}}
\newcommand{\mH}{\mathbf{H}}
% Sets and groups
\newcommand{\RR}{\mathbb{R}}
\newcommand{\SO}{\mathrm{SO}}
% Units
\newcommand{\degps}{\,\text{deg/s}}
\newcommand{\Nm}{\ensuremath{\,\text{N}\cdot\text{m}}}
\newcommand{\BW}{\,\text{BW}}

% ========== Listings ==========
\usepackage{listings}

\lstset{
  basicstyle=\ttfamily\small,
  keywordstyle=\color{defblue}\bfseries,
  commentstyle=\color{thmgreen}\itshape,
  stringstyle=\color{noteorange},
  backgroundcolor=\color{codegray},
  frame=single,
  rulecolor=\color{gray!40},
  breaklines=true,
  showstringspaces=false,
  numbers=left,
  numberstyle=\tiny\color{gray},
  tabsize=4,
  captionpos=b
}

% ========== Header/Footer ==========
\pagestyle{fancy}
\fancyhf{}
\fancyhead[L]{\leftmark}
\fancyhead[R]{\thepage}
\renewcommand{\headrulewidth}{0.4pt}

% ========== Title formatting ==========
\titleformat{\chapter}[display]
  {\normalfont\huge\bfseries\color{defblue}}
  {\chaptertitlename\ \thechapter}{20pt}{\Huge}

\titleformat{\section}
  {\normalfont\Large\bfseries\color{defblue!80!black}}
  {\thesection}{1em}{}

\titleformat{\subsection}
  {\normalfont\large\bfseries\color{defblue!60!black}}
  {\thesubsection}{1em}{}


\begin{document}

% ==================== TITLE PAGE ====================
\begin{titlepage}
\centering
\vspace*{1.5cm}
{\Large\color{gray} Part 5 --- Research Proposal}
\vspace{0.5cm}

{\Huge\bfseries\color{defblue} 3D Gaussian Splatting 기반\\[0.3cm]
마커리스 투구 동작 분석:\\[0.3cm]
마커 기반 정확도를 넘어서}

\vspace{1cm}

{\LARGE\color{defblue!70!black} GART-Pitch: Gaussian Articulated Representation\\
for High-Fidelity Markerless Pitching Biomechanics}

\vspace{1.5cm}

\begin{tikzpicture}[scale=0.55]
% Left: Marker-based (old)
\begin{scope}[shift={(-6,0)}]
  \draw[thick, gray!50] (0,0) -- (0,3) -- (-0.8,5) -- (0,4.5) -- (0.8,5) -- (0,3);
  \draw[thick, gray!50] (0,0) -- (-0.6,-3);
  \draw[thick, gray!50] (0,0) -- (0.6,-3);
  \foreach \x/\y in {0/3, -0.8/5, 0.8/5, 0/0, -0.6/-3, 0.6/-3, 0/4.5, -0.4/1.5, 0.4/1.5} {
    \fill[injuryred] (\x,\y) circle (0.12);
  }
  \node[font=\small\bfseries, text=gray] at (0,-4) {Marker-Based};
  \node[font=\tiny, text=gray] at (0,-4.7) {MPJPE $< 1$ mm};
  \node[font=\tiny, text=injuryred] at (0,-5.3) {(invasive, expensive)};
\end{scope}

% Center: Existing markerless (current SOTA)
\begin{scope}[shift={(0,0)}]
  \draw[thick, noteorange!70] (0,0) -- (0,3) -- (-0.8,5) -- (0,4.5) -- (0.8,5) -- (0,3);
  \draw[thick, noteorange!70] (0,0) -- (-0.6,-3);
  \draw[thick, noteorange!70] (0,0) -- (0.6,-3);
  \foreach \x/\y in {0/3, -0.8/5, 0.8/5, 0/0, -0.6/-3, 0.6/-3, 0/4.5} {
    \fill[noteorange] (\x,\y) circle (0.08);
  }
  \draw[noteorange, dashed, thick] (-0.8,5) circle (0.4);
  \draw[noteorange, dashed, thick] (0.8,5) circle (0.4);
  \node[font=\small\bfseries, text=noteorange!80!black] at (0,-4) {Markerless SOTA};
  \node[font=\tiny, text=noteorange!80!black] at (0,-4.7) {MPJPE 52--57 mm};
  \node[font=\tiny, text=noteorange] at (0,-5.3) {(keypoint-only)};
\end{scope}

% Right: Ours (GART-Pitch)
\begin{scope}[shift={(6,0)}]
  % Body with gaussian "glow"
  \fill[defblue!10] (0,1.5) ellipse (1.5 and 4);
  \draw[thick, defblue] (0,0) -- (0,3) -- (-0.8,5) -- (0,4.5) -- (0.8,5) -- (0,3);
  \draw[thick, defblue] (0,0) -- (-0.6,-3);
  \draw[thick, defblue] (0,0) -- (0.6,-3);
  % Dense gaussians
  \foreach \i in {1,...,30} {
    \pgfmathsetmacro{\rx}{-1.2+2.4*rand}
    \pgfmathsetmacro{\ry}{-2.5+7*rand}
    \fill[defblue!40, opacity=0.3] (\rx,\ry) circle (0.15);
  }
  \node[font=\small\bfseries, text=defblue] at (0,-4) {GART-Pitch (Ours)};
  \node[font=\tiny, text=defblue] at (0,-4.7) {Target: MPJPE $< 20$ mm};
  \node[font=\tiny, text=thmgreen] at (0,-5.3) {(dense, photometric)};
\end{scope}

% Arrows
\draw[-Stealth, very thick, thmgreen] (2.5, 0) -- (3.8, 0) node[midway, above, font=\tiny\bfseries, text=thmgreen] {$>$60\% better};
\draw[-Stealth, thick, gray, dashed] (-3.5, 0) -- (-2.2, 0);
\end{tikzpicture}

\vspace{1.5cm}
{\large 컴퓨터 비전 $\times$ 스포츠 생체역학 $\times$ 3D 재구성}
\vspace{0.5cm}

{\large 2026}

\vspace{\fill}
{\small\color{gray} 본 문서는 Part 1--4 교재의 후속 연구 제안서로,\\
3D Gaussian Splatting 기반 마커리스 투구 동작 분석이\\
기존 마커리스 SOTA를 능가하고 마커 기반 정확도에 근접할 수 있음을 논증합니다.}
\end{titlepage}

\tableofcontents
\newpage


% ==================== CHAPTER 0: 연구 요약 ====================
\chapter*{연구 요약 (Executive Summary)}
\addcontentsline{toc}{chapter}{연구 요약}

\begin{hypothesis}[title={핵심 가설}]
\textbf{7시점 240\,Hz RGB 영상에서 GART(Gaussian Articulated Representation Templates) 기반
분석-합성(analysis-by-synthesis) 자세 정제를 적용하면,
투구 동작의 3D 관절 위치 정확도(MPJPE)가 기존 마커리스 SOTA(52--57\,mm)를 60\% 이상 개선하여
20\,mm 이하를 달성하고, 임상적 핵심 지표(어깨 MER, 팔꿈치 외반 토크)에서
마커 기반(Vicon) 결과와 5$^\circ$ / 10\% 이내의 일치도를 보일 것이다.}
\end{hypothesis}

\vspace{0.5cm}

\noindent\textbf{배경.} 야구 투구는 인체에서 가장 빠른 관절 운동(어깨 내회전 7,000--9,000\degps)을 포함하며,
MLB 투수의 35.3\%가 UCL 재건술 경험이 있을 정도로 부상 위험이 높다.
정확한 생체역학 분석은 부상 예방의 핵심이지만,
현재 gold standard인 Vicon 마커 기반 시스템은 고비용, 침습적 마커 부착, 가려짐 등의 한계가 있다.
기존 마커리스 SOTA(Theia3D, KinaTrax)는 MPJPE 52--57\,mm, 회전 속도 오차 50--400\degps 이상으로
임상 의사결정에 충분하지 않다.

\vspace{0.3cm}
\noindent\textbf{방법.} 본 연구는 GART 프레임워크를 투구 동작에 특화시킨 \textbf{GART-Pitch}를 제안한다.
7대 240\,Hz 카메라 영상에서 (1) 다시점 SMPL-X 초기 피팅, (2) GART 기반 캐노니컬 Gaussian 학습,
(3) 렌더링 기반 자세 정제(photometric loss + multi-view consistency),
(4) 가상 마커 추출 및 OpenSim 역운동학/역동역학을 수행한다.
핵심 기여는 키포인트 재투영만이 아닌 \textbf{전체 픽셀 수준의 외형 일치}를 감독 신호로 활용하여
자세 정확도를 비약적으로 향상시키는 것이다.

\vspace{0.3cm}
\noindent\textbf{검증 계획.}
Vicon 마커 기반 역운동학을 gold standard로 하여 MPJPE, PA-MPJPE, 관절각 RMSE,
Bland--Altman 분석, ICC를 보고한다.
기존 SOTA(Theia3D, KinaTrax, OpenPose+삼각측량, MotionBERT)와 직접 비교하며,
특히 투구의 가속 단계($\sim$30\,ms)에서의 정확도를 집중 평가한다.

\vspace{0.3cm}
\noindent\textbf{기대 영향.}
마커 없이 Vicon에 근접하는 정확도로 투구 분석이 가능해지면,
(1) 경기 중 실시간 분석, (2) 과거 영상 소급 분석, (3) 비용 대폭 절감,
(4) 선수의 자연스러운 동작 보장이라는 네 가지 혁신이 실현된다.


% ====================================================================
% CHAPTER 1: 서론
% ====================================================================
\chapter{서론: 문제 정의와 연구 동기}
\label{ch:introduction}

\section{투구 생체역학 분석의 현황과 한계}

야구 투구는 인체 관절 운동 중 가장 극단적인 움직임을 포함한다.
어깨 내회전 각속도는 7,000--9,000\degps에 달하며,
팔꿈치에 작용하는 외반 토크는 체중의 50\% 이상에 해당하는 64\Nm에 이른다.
이 극단적 부하의 정확한 측정은 부상 예방 전략 수립의 전제 조건이다.

현재 생체역학 분석의 gold standard는 \textbf{Vicon 마커 기반 모션캡처}로,
반사 마커를 신체에 부착하고 적외선 카메라로 추적하여 $< 1$\,mm의 관절 위치 정확도를 달성한다.
그러나 이 방법은 다음과 같은 근본적 한계를 갖는다:

\begin{itemize}[leftmargin=*]
  \item \textbf{고비용}: 시스템 가격 \$100,000+, 전문 인력 필요
  \item \textbf{침습적 마커}: 피부 위 마커가 고유수용감각에 영향을 줄 수 있음
  \item \textbf{마커 가려짐}: 와인드업, 팔로스루에서 마커가 가려지거나 탈락
  \item \textbf{연구실 한정}: 경기장에서의 실시간 분석 불가능
  \item \textbf{연조직 아티팩트(STA)}: 빠른 움직임 시 피부 위 마커가 뼈와 다르게 움직임
\end{itemize}

\section{기존 마커리스 접근법의 한계}
\label{sec:existing_markerless}

마커리스 시스템은 이러한 한계를 극복하고자 등장했으나, 투구 분석에서는 아직 충분한 정확도에 도달하지 못했다.

\begin{table}[H]
\centering
\caption{기존 마커리스 시스템의 투구 분석 성능 (2025년 검증 연구 기준)}
\label{tab:existing_sota}
\renewcommand{\arraystretch}{1.3}
\small
\begin{tabular}{p{3.5cm}cccc}
\toprule
\textbf{시스템} & \textbf{MPJPE (mm)} & \textbf{관절각 RMSE} & \textbf{회전속도 오차} & \textbf{방식} \\
\midrule
Vicon (gold std.) & $< 1$ & 기준 & 기준 & 마커 기반 \\
Theia3D & 52.0 $\pm$ 12.3 & 2--16$^\circ$ & 50--400+\degps & 마커리스 \\
KinaTrax & 56.6 $\pm$ 9.4 & 2--16$^\circ$ & 50--400+\degps & 마커리스 \\
OpenPose + 삼각측량 & $\sim$70--90 & 5--25$^\circ$ & 검증 없음 & 키포인트 \\
MediaPipe & $\sim$80--100 & 8--30$^\circ$ & 검증 없음 & 키포인트 \\
\bottomrule
\end{tabular}
\end{table}

\begin{clinicalnote}
마커리스 시스템의 MPJPE 52--57\,mm는 보행 분석에서는 수용 가능하지만,
투구 분석에서는 임상적으로 중요한 한계이다.
특히 \textbf{팔꿈치 외반 토크} 계산은 모멘트 팔(moment arm)에 매우 민감하여,
수 mm의 관절 중심 오차가 수 Nm의 토크 오차로 증폭된다.
또한 투구에서 가장 중요한 \textbf{어깨 회전 변수}(MER, 내회전 속도)에서
가장 큰 불일치가 관찰된다.
\end{clinicalnote}

\subsection{기존 마커리스의 근본적 병목}

기존 마커리스 시스템의 한계는 \textbf{감독 신호의 빈곤}에서 기인한다:

\begin{enumerate}[leftmargin=*]
  \item \textbf{희소 키포인트 의존}: OpenPose/MediaPipe는 17--33개 키포인트만 추출.
        7,000\degps 어깨 회전에서 키포인트 2--3개로 회전각을 결정하는 것은 본질적으로 부정확.
  \item \textbf{모션 블러 취약}: 30\,fps용으로 설계된 검출기를 240\,fps에 적용해도,
        가속 단계의 팔은 흐릿한 줄무늬로 나타남.
  \item \textbf{깊이 모호성}: 단안 또는 적은 수의 뷰에서 깊이 추정이 근본적으로 어려움.
  \item \textbf{외형 정보 미활용}: 키포인트 좌표만 사용하고, 이미지의 풍부한 외형 정보
        (근육 윤곽, 유니폼 주름, 피부 텍스처)를 활용하지 않음.
\end{enumerate}


\section{제안하는 접근법: Analysis-by-Synthesis}
\label{sec:proposed_approach}

\begin{keyidea}
본 연구의 핵심 통찰은 다음과 같다:
\textbf{3DGS/GART 기반 렌더링이 제공하는 전체 픽셀 수준의 광도 손실(photometric loss)은
희소 키포인트 재투영 손실보다 본질적으로 더 풍부한 감독 신호를 제공한다.}

\vspace{0.3cm}
\noindent $1920 \times 1080$ 해상도 $\times$ 7대 카메라 = 프레임당 약 \textbf{14.5M 픽셀}의 감독 신호.\\
키포인트 기반: 33개 관절 $\times$ 7대 카메라 = 프레임당 \textbf{231개} 좌표값.\\[0.2cm]
\textbf{감독 신호 밀도 차이: $\sim$63,000$\times$}
\end{keyidea}

이 감독 신호의 비약적 차이가 자세 정확도의 비약적 향상을 가능하게 한다는 것이 본 연구의 가설이다.

\begin{figure}[H]
\centering
\begin{tikzpicture}[
  block/.style={rectangle, rounded corners, draw, minimum width=3cm, minimum height=0.8cm, align=center, font=\small},
  arrow/.style={-{Stealth[length=2.5mm]}, thick}
]
% Keypoint-based (top row)
\node[block, fill=noteorange!20, draw=noteorange] (kp_det) at (0,2) {2D Keypoint\\Detection};
\node[block, fill=noteorange!20, draw=noteorange] (kp_tri) at (4.5,2) {Multi-view\\Triangulation};
\node[block, fill=noteorange!20, draw=noteorange] (kp_smpl) at (9,2) {SMPL-X\\Fitting};
\node[font=\small\bfseries, text=noteorange!80!black] at (-3,2) {기존 방법};
\node[font=\tiny, text=gray] at (9, 1.2) {감독 신호: 231개 좌표};

\draw[arrow, noteorange] (kp_det) -- (kp_tri);
\draw[arrow, noteorange] (kp_tri) -- (kp_smpl);

% GART-based (bottom row)
\node[block, fill=defblue!15, draw=defblue] (g_init) at (0,-1) {초기 SMPL-X\\피팅};
\node[block, fill=defblue!15, draw=defblue] (g_gart) at (4.5,-1) {GART 학습\\+ Rendering};
\node[block, fill=defblue!15, draw=defblue] (g_ref) at (9,-1) {Photometric\\Pose Refinement};
\node[font=\small\bfseries, text=defblue] at (-3,-1) {GART-Pitch};
\node[font=\tiny, text=gray] at (9, -1.8) {감독 신호: 14.5M 픽셀};

\draw[arrow, defblue] (g_init) -- (g_gart);
\draw[arrow, defblue] (g_gart) -- (g_ref);

% Feedback loop
\draw[arrow, darkred, dashed] (g_ref.south) -- ++(0,-0.5) -| (g_gart.south)
  node[midway, below, font=\tiny\color{darkred}] {iterative refinement};

% Comparison
\draw[very thick, thmgreen, -Stealth] (11, 2) -- (11, -1)
  node[midway, right, font=\small\bfseries\color{thmgreen}] {$63,\!000\times$\\감독 신호};
\end{tikzpicture}
\caption{기존 키포인트 기반 방법 vs GART-Pitch의 감독 신호 밀도 비교.
GART-Pitch는 프레임당 14.5M 픽셀의 광도 일치를 요구하여 자세 정제의 정확도를 비약적으로 향상시킨다.}
\label{fig:supervision_comparison}
\end{figure}


\section{연구 기여 (Contributions)}

본 연구의 구체적 기여는 다음과 같다:

\begin{enumerate}[leftmargin=*]
  \item \textbf{GART-Pitch 프레임워크}: 3D Gaussian Splatting 기반 인체 재구성을
        고속 투구 동작 분석에 최초로 적용하는 end-to-end 파이프라인 설계.

  \item \textbf{시간적 정규화 + 적응형 자세 정제}: 투구의 6단계 특성(가속 30\,ms)에
        맞춘 위상별 적응형 시간 정규화와 coarse-to-fine 자세 정제 전략.

  \item \textbf{포괄적 벤치마크}: Vicon gold standard 대비, 기존 마커리스 SOTA 4종
        (Theia3D, KinaTrax, OpenPose+삼각측량, MotionBERT)과의 체계적 정량 비교.

  \item \textbf{임상 지표 검증}: MPJPE를 넘어 투구 특화 임상 지표
        (MER, 팔꿈치 외반 토크, 어깨 내회전 속도)에서의 정확도 검증.
\end{enumerate}


% ====================================================================
% CHAPTER 2: 관련 연구
% ====================================================================
\chapter{관련 연구}
\label{ch:related_work}

\section{투구 생체역학 분석의 발전}

투구 생체역학 연구는 세 세대에 걸쳐 발전해왔다.

\subsection{제1세대: 마커 기반 모션캡처 (1990s--현재)}

Fleisig 등(1996)이 확립한 Vicon 기반 분석은 여전히 gold standard로 사용된다.
ASMI(American Sports Medicine Institute)의 프로토콜은 대부분의 투구 생체역학 연구의 기초이며,
UCL 토크 $> 64$\Nm이 부상 위험 요인이라는 핵심 발견도 이 세대에서 이루어졌다.
%% [CITATION: Fleisig1996 -- PLACEHOLDER, VERIFY]

\subsection{제2세대: 상용 마커리스 시스템 (2018--현재)}

KinaTrax(8대 300\,fps, MLB 경기장 설치)와 Theia3D(다시점 RGB 기반)가 대표적이다.
2025년 검증 연구에서 18명 대학 투수를 대상으로 Vicon과 비교한 결과,
MPJPE 52--57\,mm, 관절각 RMSE 2--16$^\circ$로 보고되었다.
어깨 회전 변수에서 가장 큰 불일치가 관찰되며,
이는 정확히 투구 분석에서 가장 중요한 변수(MER, 내회전 속도)이다.
%% [CITATION: Boddy2025_KinaTrax_validation -- PLACEHOLDER, VERIFY]

\subsection{제3세대: 딥러닝 기반 마커리스 (2020--현재)}

단안 영상 기반 인체 자세 추정(HMR, PARE, 4DHumans, MotionBERT)이 급속히 발전했으나,
투구와 같은 극한 동작에서의 검증은 제한적이다.
주목할 만한 연구로, Bright 등(2026)이 MLB 방송 영상(단안, 30\,fps)만으로
18개 생체역학 지표를 추출하여 토미 존 수술 예측 AUC 0.811을 달성했다.
이는 저품질 입력에서도 의미 있는 분석이 가능함을 시사하지만,
지표 수준의 정확도(개별 관절 각도)는 여전히 Vicon에 미치지 못한다.
%% [CITATION: Bright2026_broadcast -- PLACEHOLDER, VERIFY]


\section{3D Gaussian Splatting과 인체 재구성}

\subsection{3DGS 기초}

Kerbl 등(2023)이 제안한 3D Gaussian Splatting은 Neural Radiance Fields(NeRF)의 암시적 표현을
명시적 3D Gaussian의 집합으로 대체하여 실시간 렌더링과 빠른 학습을 가능하게 했다.
각 Gaussian $g_k = \{\vmu_k, \mSigma_k, \alpha_k, \vc_k\}$는 중심, 공분산, 불투명도, 색상(SH 계수)으로 정의된다.
%% [CITATION: Kerbl2023_3DGS -- PLACEHOLDER, VERIFY]

\subsection{동적 인체 재구성: GART와 관련 연구}

\textbf{GART}(Lei 등, 2024)는 SMPL의 Linear Blend Skinning(LBS) 가중치를 활용하여
정규 공간의 Gaussian을 자세별로 변환하는 구조를 제안했다.
GauHuman(Hu 등, 2024), ExAvatar(Moon 등, 2024) 등이 유사한 접근을 발전시켰으며,
모두 photometric loss 기반 자세 정제를 지원한다.
%% [CITATION: Lei2024_GART, Hu2024_GauHuman, Moon2024_ExAvatar -- PLACEHOLDER, VERIFY]

그러나 기존 연구들은 대부분 \textbf{보행이나 댄스 등 상대적으로 느린 동작}에서 검증되었으며,
7,000\degps 이상의 고속 관절 운동에서의 성능은 검증된 바 없다.
이것이 본 연구가 채우고자 하는 핵심 갭이다.


\section{감독 신호 밀도와 자세 정확도}

자세 추정의 감독 신호 밀도가 정확도에 미치는 영향은 이론적으로나 실증적으로 지지된다:

\begin{itemize}[leftmargin=*]
  \item \textbf{키포인트 기반}: $\sim$20--33개 관절 $\times$ 뷰 수. 자유도 대비 관측 수가 부족하여
        ill-posed 문제가 되기 쉬움.
  \item \textbf{실루엣 기반}: $\sim$10K--50K 경계 픽셀. 외형의 내부 구조는 활용 불가.
  \item \textbf{광도 기반 (3DGS)}: $\sim$수백만 픽셀. 표면 텍스처, 근육 윤곽, 그림자 등
        모든 시각 정보를 감독 신호로 활용.
\end{itemize}

\begin{table}[H]
\centering
\caption{감독 신호 유형별 밀도 비교 (7대 1080p 카메라 기준)}
\label{tab:supervision_density}
\begin{tabular}{lccc}
\toprule
\textbf{감독 유형} & \textbf{프레임당 신호 수} & \textbf{차원} & \textbf{대표 시스템} \\
\midrule
키포인트 좌표 & $\sim$231 & 2D & OpenPose, Theia3D \\
실루엣 경계 & $\sim$70K & 1D (binary) & Theia3D \\
\textbf{전체 광도 (Ours)} & \textbf{$\sim$14.5M} & \textbf{3D (RGB)} & \textbf{GART-Pitch} \\
\bottomrule
\end{tabular}
\end{table}


% ====================================================================
% CHAPTER 3: 방법론
% ====================================================================
\chapter{GART-Pitch: 제안 방법론}
\label{ch:method}

\section{전체 파이프라인 개요}

\begin{figure}[H]
\centering
\begin{tikzpicture}[
  block/.style={rectangle, draw, rounded corners=3pt, minimum width=2.8cm, minimum height=0.8cm,
    align=center, font=\small},
  data/.style={rectangle, draw, dashed, rounded corners=3pt, minimum width=2.2cm, minimum height=0.6cm,
    align=center, font=\scriptsize},
  arrow/.style={-{Stealth[length=2.5mm]}, thick},
  node distance=0.6cm and 1.5cm
]
% Row 1: Inputs
\node[block, fill=lightorange] (rgb) at (0,0) {7$\times$ RGB\\240\,Hz};
\node[block, fill=lightgreen] (vicon) at (6,0) {Vicon\\250--500\,Hz};

% Row 2: Preprocessing
\node[block, fill=lightgray] (sync) at (3,-1.5) {시간/공간\\동기화};
\node[block, fill=lightblue] (pose2d) at (0,-3) {Multi-view\\2D Pose (ViTPose)};
\node[block, fill=lightblue] (calib) at (6,-3) {카메라 보정\\(COLMAP)};

% Row 3: SMPL-X Initialization
\node[block, fill=cvpurple!15, draw=cvpurple] (smpl_init) at (3,-4.5) {다시점 SMPL-X\\초기 피팅};

% Row 4: GART Core
\node[block, fill=defblue!20, draw=defblue, minimum width=8cm] (gart) at (3,-6.5)
  {\textbf{GART-Pitch Core}\\
   Canonical Gaussian $\xrightarrow{\text{LBS}}$ Posed Gaussian $\xrightarrow{\text{Splat}}$ Rendered Image};

% Row 5: Pose Refinement
\node[block, fill=darkred!15, draw=darkred] (refine) at (3,-8.5)
  {\textbf{위상별 적응형 자세 정제}\\
   $\vtheta^* = \arg\min \mathcal{L}_{\text{photo}} + \mathcal{L}_{\text{temporal}} + \mathcal{L}_{\text{biomech}}$};

% Row 6: Biomechanics
\node[block, fill=biogreen!20, draw=biogreen] (bio) at (-1,-10.5) {ISB 관절각\\+ OpenSim IK/ID};
\node[block, fill=biogreen!20, draw=biogreen] (clinical) at (3,-12) {투구 임상 지표\\MER, $\tau_{\text{valgus}}$, $\omega_{\text{IR}}$};
\node[block, fill=thmgreen!20, draw=thmgreen] (validate) at (7,-10.5) {Vicon 대비\\정량적 검증};

% Arrows
\draw[arrow] (rgb) -- (sync);
\draw[arrow] (vicon) -- (sync);
\draw[arrow] (sync) -- (pose2d);
\draw[arrow] (sync) -- (calib);
\draw[arrow] (pose2d) -- (smpl_init);
\draw[arrow] (calib) -- (smpl_init);
\draw[arrow] (smpl_init) -- (gart);
\draw[arrow] (gart) -- (refine);
\draw[arrow, darkred, dashed] (refine.east) -- ++(1.5,0) |- (gart.east)
  node[pos=0.25, right, font=\tiny\color{darkred}] {iterative};
\draw[arrow] (refine) -- (bio);
\draw[arrow] (bio) -- (clinical);
\draw[arrow] (refine) -- (validate);
\draw[arrow] (vicon.south) -- ++(0,-8) -- (validate.north);

\end{tikzpicture}
\caption{GART-Pitch 전체 파이프라인. 핵심은 GART Core의 렌더링 기반 자세 정제(빨간 점선 루프)이다.}
\label{fig:full_pipeline}
\end{figure}


\section{Stage 1: 데이터 전처리 및 초기화}

\subsection{다시점 동기화}

7대 240\,Hz 카메라(6대 landscape 1920$\times$1080, 1대 portrait 1080$\times$1920)와
Vicon(250--500\,Hz)의 시간 동기화는 오디오 기반 크로스-상관 방법으로 수행한다:
\begin{equation}
  \Delta t^* = \arg\max_{\Delta t} \sum_t s_{\text{audio}}^{(i)}(t) \cdot s_{\text{audio}}^{(j)}(t + \Delta t)
\end{equation}
공간 동기화는 체커보드 보정 + DLT(Direct Linear Transform)로 Vicon 좌표계와 RGB 좌표계를 정합한다.

\subsection{카메라 보정}

COLMAP을 이용한 SfM(Structure from Motion)으로 내부/외부 파라미터를 추정한다.
투구 세팅에서는 피험자가 없는 빈 장면을 촬영하여 충분한 특징점을 확보하는 것이 중요하다.
재투영 오차 $< 0.5$\,px를 목표로 한다.

\subsection{초기 SMPL-X 피팅}

ViTPose로 각 카메라의 2D 관절을 추출한 후, EasyMocap 프레임워크로 다시점 SMPL-X 피팅을 수행한다:
\begin{equation}
(\vtheta^{(0)}, \vbeta^*) = \arg\min_{\vtheta, \vbeta}
  \underbrace{\sum_{c=1}^{7}\sum_{j=1}^{J} w_{cj}\|\pi_c(J_j(\vtheta,\vbeta)) - \hat{\vx}_{cj}\|^2}_{\text{2D 재투영}}
  + \lambda_\beta\|\vbeta\|^2 + \lambda_\theta\|\vtheta\|^2
\label{eq:init_smpl}
\end{equation}

이 초기 피팅은 GART 학습의 시작점으로만 사용되며,
최종 정확도는 이후 렌더링 기반 정제 단계에서 비약적으로 향상된다.


\section{Stage 2: GART 캐노니컬 Gaussian 학습}

\subsection{캐노니컬 공간 Gaussian}

정규(canonical) 공간에서 $N$ 개의 3D Gaussian을 학습한다:
\begin{equation}
  \mathcal{G}_{\text{can}} = \{g_k\}_{k=1}^{N}, \quad
  g_k = \{\vmu_k^{\text{can}}, \mSigma_k^{\text{can}}, \alpha_k, \vc_k\}
\end{equation}
초기 Gaussian은 SMPL-X 메시 정점에서 샘플링하여 초기화한다 ($N_0 \approx 10,475$, SMPL-X 정점 수).

\subsection{LBS 기반 자세 변환}

프레임 $t$의 자세 $\vtheta_t$에 대해, Gaussian 중심과 공분산을 변환한다:
\begin{align}
  \vmu_k^{\text{posed}}(t) &= \Big(\sum_{b=1}^{B} w_{kb}\, \mT_b(\vtheta_t)\Big) \bar{\vmu}_k^{\text{can}} \\
  \mSigma_k^{\text{posed}}(t) &= \mR_{\text{lbs},k}(t)\, \mSigma_k^{\text{can}}\, \mR_{\text{lbs},k}(t)^\top
\end{align}
여기서 $w_{kb}$는 Gaussian $k$의 본 $b$에 대한 스키닝 가중치 (SMPL-X 기본값 + 학습 가능 오프셋 $\Delta w_{kb}$),
$\mT_b \in SE(3)$는 본 변환 행렬, $\bar{\vmu}$는 동차좌표 확장이다.

\subsection{미분 가능 래스터화}

변환된 Gaussian을 미분 가능하게 래스터화하여 각 카메라 뷰의 이미지를 생성한다:
\begin{equation}
  \hat{I}_c(t) = \text{Splat}(\mathcal{G}_{\text{posed}}(t),\; \mK_c, [\mR_c | \vt_c])
\end{equation}
여기서 $\mK_c$는 카메라 $c$의 내부 행렬, $[\mR_c|\vt_c]$는 외부 행렬이다.


\section{Stage 3: 위상별 적응형 자세 정제}
\label{sec:phase_adaptive}

이것이 GART-Pitch의 핵심 기여이다.
투구 동작의 각 단계별로 다른 정규화 전략을 적용하여 정확도를 극대화한다.

\subsection{투구 위상 자동 검출}

투구의 6단계(와인드업, 스트라이드, 후기 코킹, 가속, 감속, 팔로스루)를
자동으로 분할하기 위해, 초기 SMPL-X 피팅의 관절 속도 프로파일을 활용한다:

\begin{itemize}[nosep]
  \item \textbf{발 접지(FC)}: 앞발 관절의 수직 속도가 0을 교차하는 시점
  \item \textbf{MER}: 어깨 외회전 각도의 최대값
  \item \textbf{공 릴리스(BR)}: 손목 관절의 전방 속도 최대값 직후
\end{itemize}

\subsection{위상별 적응형 손실 함수}

\begin{equation}
\mathcal{L}_{\text{total}}(t) =
  \underbrace{\mathcal{L}_{\text{photo}}(t)}_{\text{광도 일치}}
  + \underbrace{\lambda_{\text{temp}}^{\phi(t)} \mathcal{L}_{\text{temp}}(t)}_{\text{위상별 시간 정규화}}
  + \underbrace{\lambda_{\text{bio}} \mathcal{L}_{\text{bio}}(t)}_{\text{생체역학 제약}}
\label{eq:total_loss}
\end{equation}

\subsubsection{광도 손실 (Photometric Loss)}

\begin{equation}
\mathcal{L}_{\text{photo}}(t) = \sum_{c=1}^{7}
  \Big[ (1-\lambda_s)\|I_c(t) - \hat{I}_c(t)\|_1 + \lambda_s\, \mathcal{L}_{\text{D-SSIM}}(I_c(t), \hat{I}_c(t)) \Big]
\end{equation}

7대 카메라 모두의 광도 일치를 동시에 요구하므로,
단일 뷰에서 모호한 자세도 다시점 일관성으로 해소된다.

\subsubsection{위상별 시간 정규화 (Phase-Adaptive Temporal Regularization)}

\begin{equation}
\mathcal{L}_{\text{temp}}(t) = \|\vtheta_t - \vtheta_{t-1}\|^2 + \gamma \|\vtheta_t - 2\vtheta_{t-1} + \vtheta_{t-2}\|^2
\end{equation}

핵심은 위상별로 $\lambda_{\text{temp}}^{\phi(t)}$를 적응적으로 조절하는 것이다:

\begin{table}[H]
\centering
\caption{투구 위상별 시간 정규화 가중치}
\label{tab:phase_weights}
\begin{tabular}{lcl}
\toprule
\textbf{위상 $\phi(t)$} & \textbf{$\lambda_{\text{temp}}^{\phi}$} & \textbf{근거} \\
\midrule
와인드업 & 10.0 & 느린 동작, 강한 정규화 가능 \\
스트라이드 & 5.0 & 중간 속도 \\
후기 코킹 & 2.0 & 빠른 외회전 시작, 정규화 완화 \\
\textbf{가속} & \textbf{0.5} & \textbf{7,000\degps, 급격한 변화 허용 필수} \\
감속 & 1.0 & 급격한 속도 변화, 중간 정규화 \\
팔로스루 & 3.0 & 점진적 감속 \\
\bottomrule
\end{tabular}
\end{table}

\begin{importantnote}
가속 단계(약 30\,ms, 7프레임)에서 시간 정규화를 과도하게 적용하면,
7,000\degps의 실제 각속도 변화를 억제하여 오히려 정확도가 크게 저하된다.
이것이 기존 방법들이 어깨 회전 변수에서 가장 큰 오차를 보이는 주요 원인 중 하나이다.
GART-Pitch는 이 구간에서 정규화를 최소화하고 광도 손실에 의존한다.
\end{importantnote}

\subsubsection{생체역학적 제약 (Biomechanical Constraints)}

해부학적으로 불가능한 관절 운동을 방지하는 소프트 제약:
\begin{equation}
\mathcal{L}_{\text{bio}}(t) =
  \sum_{k \in \mathcal{J}} \max(0, \theta_k(t) - \theta_k^{\max})^2 + \max(0, \theta_k^{\min} - \theta_k(t))^2
\end{equation}
여기서 $\theta_k^{\max}, \theta_k^{\min}$은 관절 $k$의 해부학적 가동범위이다.
예: 팔꿈치 과신전 $< 5^\circ$, 어깨 외전 $< 180^\circ$.

\subsection{Coarse-to-Fine 정제 전략}

자세 정제는 두 단계로 수행한다:

\begin{enumerate}[leftmargin=*]
  \item \textbf{Phase 1 (Coarse)}: Gaussian 파라미터 고정, 자세 $\vtheta_t$만 최적화.
        큰 학습률($\eta = 10^{-3}$)로 전역적 자세 오차를 빠르게 교정. 500 iterations.

  \item \textbf{Phase 2 (Fine)}: Gaussian 파라미터와 자세를 공동 최적화.
        작은 자세 학습률($\eta = 10^{-5}$)로 세밀한 정합.
        Gaussian 학습률은 일반적 3DGS 수준 유지. 3000 iterations.
\end{enumerate}


\section{Stage 4: 생체역학 지표 산출}

정제된 SMPL-X 자세 $\vtheta^*$에서 투구 생체역학 지표를 산출한다.

\subsection{관절 각도 추출 (ISB 분해)}

어깨 관절의 $Y$-$X'$-$Y''$ 오일러 분해:
\begin{equation}
  \mR_{\text{shoulder}} = \mR_y(\alpha)\,\mR_{x'}(\beta)\,\mR_{y''}(\gamma)
\end{equation}
팔꿈치의 $Z$-$X$-$Y$ 분해:
\begin{equation}
  \mR_{\text{elbow}} = \mR_z(\text{flex/ext})\,\mR_x(\text{carry})\,\mR_y(\text{pro/sup})
\end{equation}

\subsection{핵심 투구 지표}

\begin{table}[H]
\centering
\caption{산출 대상 투구 생체역학 지표}
\label{tab:pitching_metrics}
\renewcommand{\arraystretch}{1.3}
\begin{tabular}{p{4cm}lcc}
\toprule
\textbf{지표} & \textbf{기호} & \textbf{정상 범위} & \textbf{부상 임계값} \\
\midrule
어깨 최대 외회전 & MER & 160--180$^\circ$ & $> 185^\circ$ \\
어깨 내회전 속도 & $\omega_{\text{IR}}$ & 7,000--9,000\degps & --- \\
팔꿈치 외반 토크 & $\tau_{\text{valgus}}$ & 50--64\Nm & $> 64$\Nm \\
어깨 내회전 토크 & $\tau_{\text{IR}}$ & 50--80\Nm & --- \\
스트라이드 길이 & $L_s$ & 85\% 신장 & $< 74\%$ \\
몸통 분리 & $\Delta\theta_{\text{trunk}}$ & 40--60$^\circ$ & $< 30^\circ$ \\
\bottomrule
\end{tabular}
\end{table}

\subsection{가상 마커 + OpenSim 파이프라인}

SMPL-X 메시에서 해부학적 랜드마크에 해당하는 정점을 추출하여 TRC 형식으로 변환한 후,
OpenSim의 역운동학(IK)과 역동역학(ID)을 수행하여 관절 토크를 산출한다.
이를 통해 마커리스 파이프라인의 결과를 전통적 생체역학 분석 도구로 직접 검증할 수 있다.


% ====================================================================
% CHAPTER 4: 실험 설계
% ====================================================================
\chapter{실험 설계}
\label{ch:experiments}

\section{데이터 수집}

\begin{setupbox}[title={실험 세팅}]
\begin{itemize}[leftmargin=*]
  \item \textbf{피험자}: 대학/프로 수준 투수 15--20명 (IRB 승인 후)
  \item \textbf{RGB 카메라}: 7대, 240\,Hz
  \begin{itemize}
    \item 6대: 1920$\times$1080 (landscape)
    \item 1대: 1080$\times$1920 (portrait, 투수 후방)
  \end{itemize}
  \item \textbf{Vicon}: 12+대, 250--500\,Hz, 반사 마커 39개 (Upper Extremity Marker Set)
  \item \textbf{동기화}: 외부 트리거 + 오디오 기반 크로스-상관
  \item \textbf{투구}: 피험자당 최대 노력 패스트볼 10구 이상
  \item \textbf{환경}: 실내 투구 연구실, 충분한 조명 (마운드 + 홈플레이트)
\end{itemize}
\end{setupbox}


\section{비교 대상 (Baselines)}

\begin{table}[H]
\centering
\caption{비교 실험 설계}
\label{tab:baselines}
\renewcommand{\arraystretch}{1.3}
\begin{tabular}{p{3.5cm}p{3cm}p{6cm}}
\toprule
\textbf{방법} & \textbf{유형} & \textbf{설명} \\
\midrule
Vicon + IK & Gold Standard & 반사 마커 기반 역운동학 (모든 비교의 기준) \\
\midrule
Theia3D & Markerless SOTA & 상용 마커리스 시스템, 동일 영상 입력 \\
OpenPose + 삼각측량 & Keypoint & 키포인트 추출 후 다시점 삼각측량 \\
ViTPose + EasyMocap & Keypoint+SMPL & 다시점 SMPL-X 피팅 (우리의 초기화) \\
MotionBERT & Lifting & 단안 3D 자세 추정 후 다시점 앙상블 \\
\textbf{GART-Pitch (Ours)} & \textbf{3DGS} & \textbf{GART + 위상별 자세 정제} \\
\bottomrule
\end{tabular}
\end{table}


\section{평가 지표}

\subsection{3D 자세 정확도}

\begin{align}
  \text{MPJPE} &= \frac{1}{TJ}\sum_{t,j}\|\hat{\vp}_j(t) - \vp_j^{\text{gt}}(t)\|_2 \label{eq:mpjpe} \\
  \text{PA-MPJPE} &= \frac{1}{TJ}\sum_{t,j}\|s^*\mR^*\hat{\vp}_j(t) + \vt^* - \vp_j^{\text{gt}}(t)\|_2 \label{eq:pampjpe}
\end{align}

\subsection{관절각 정확도}

\begin{equation}
  \text{RMSE}_{\theta_k} = \sqrt{\frac{1}{T}\sum_{t=1}^{T}(\theta_k^{\text{pred}}(t) - \theta_k^{\text{gt}}(t))^2}
\end{equation}

\subsection{운동역학 정확도}

\begin{equation}
  \text{nRMSE}_\tau = \frac{\text{RMSE}_\tau}{\tau_{\max}^{\text{gt}} - \tau_{\min}^{\text{gt}}} \times 100\%
\end{equation}

\subsection{일치도 분석}

\begin{itemize}[nosep]
  \item \textbf{Bland--Altman 분석}: 바이어스 + 95\% 일치한계(LoA)
  \item \textbf{ICC(3,1)}: $> 0.90$ (우수), $0.75$--$0.90$ (양호)
  \item \textbf{CCC (Concordance Correlation Coefficient)}: 일치도의 엄격한 평가
\end{itemize}


\section{실험 프로토콜}

\subsection{실험 1: 전체 정확도 비교}

모든 방법의 전체 투구 동작에 대한 MPJPE, PA-MPJPE, 관절각 RMSE를 비교한다.

\begin{hypothesis}[title={실험 1 가설}]
GART-Pitch는 MPJPE $< 20$\,mm를 달성하여 Theia3D(52\,mm) 대비 60\% 이상 개선된다.
\end{hypothesis}

\subsection{실험 2: 투구 위상별 정확도 분석}

투구의 6단계 각각에서 관절각 RMSE를 보고한다.
특히 가속 단계(30\,ms)에서의 정확도를 집중 분석한다.

\begin{hypothesis}[title={실험 2 가설}]
가속 단계에서 GART-Pitch와 기존 SOTA의 정확도 차이가 가장 크다.
기존 방법들이 가장 실패하는 구간에서 광도 기반 정제가 가장 큰 효과를 발휘한다.
\end{hypothesis}

\subsection{실험 3: 임상 지표 일치도}

MER, 팔꿈치 외반 토크, 어깨 내회전 속도에 대해 Vicon 대비 Bland--Altman 분석을 수행한다.

\begin{hypothesis}[title={실험 3 가설}]
MER에서 RMSE $< 5^\circ$, 외반 토크에서 nRMSE $< 10\%$를 달성하여
임상 의사결정에 충분한 정확도를 보인다.
\end{hypothesis}

\subsection{실험 4: Ablation Study}

각 구성 요소의 기여도를 분석한다:

\begin{table}[H]
\centering
\caption{Ablation 실험 설계}
\label{tab:ablation}
\begin{tabular}{lccc}
\toprule
\textbf{변형} & \textbf{위상별 $\lambda_{\text{temp}}$} & \textbf{$\mathcal{L}_{\text{bio}}$} & \textbf{Coarse-to-Fine} \\
\midrule
Full GART-Pitch & \checkmark & \checkmark & \checkmark \\
w/o Phase-Adaptive & $\times$ (상수) & \checkmark & \checkmark \\
w/o Biomech. Constraint & \checkmark & $\times$ & \checkmark \\
w/o Coarse-to-Fine & \checkmark & \checkmark & $\times$ (single-stage) \\
SMPL-X Only (초기화) & $\times$ & $\times$ & $\times$ \\
\bottomrule
\end{tabular}
\end{table}


% ====================================================================
% CHAPTER 5: 예상 결과
% ====================================================================
\chapter{예상 결과 및 분석}
\label{ch:expected_results}

\section{실험 1: 전체 정확도}

\begin{resultbox}[title={예상 결과: 3D 자세 정확도}]
\begin{center}
\renewcommand{\arraystretch}{1.3}
\begin{tabular}{lccc}
\toprule
\textbf{방법} & \textbf{MPJPE (mm) $\downarrow$} & \textbf{PA-MPJPE (mm) $\downarrow$} & \textbf{관절각 RMSE $\downarrow$} \\
\midrule
OpenPose + 삼각측량 & 75--90 & 55--65 & 8--25$^\circ$ \\
MotionBERT 앙상블 & 60--75 & 45--55 & 6--20$^\circ$ \\
Theia3D & 52 $\pm$ 12 & 40 $\pm$ 8 & 2--16$^\circ$ \\
ViTPose + EasyMocap & 40--50 & 30--40 & 3--12$^\circ$ \\
\textbf{GART-Pitch (Ours)} & \textbf{15--20} & \textbf{10--15} & \textbf{1--5$^\circ$} \\
\midrule
Vicon (Gold Standard) & $< 1$ & $< 1$ & 기준 \\
\bottomrule
\end{tabular}
\end{center}
\end{resultbox}

\begin{keyidea}
GART-Pitch가 15--20\,mm MPJPE를 달성하면,
기존 마커리스 SOTA 대비 \textbf{60--70\% 개선}이며,
Vicon과의 격차를 \textbf{52\,mm $\rightarrow$ 20\,mm으로 60\% 좁히는} 것이다.
이 수준은 투구 분석에서 ``수용 가능'' 기준(40\,mm)을 넘어 ``우수'' 기준(20\,mm)에 도달한다.
\end{keyidea}


\section{실험 2: 위상별 정확도}

\begin{figure}[H]
\centering
\begin{tikzpicture}
\begin{axis}[
  width=14cm, height=7cm,
  ybar,
  bar width=0.25cm,
  xlabel={투구 위상},
  ylabel={어깨 관절각 RMSE ($^\circ$)},
  ymin=0, ymax=30,
  symbolic x coords={Wind-up, Stride, Late Cocking, Acceleration, Deceleration, Follow-through},
  xtick=data,
  x tick label style={rotate=30, anchor=east, font=\small},
  legend style={at={(0.5,-0.35)}, anchor=north, legend columns=3, font=\small},
  grid=major,
  grid style={gray!20},
  title={\small 투구 위상별 어깨 관절각 RMSE (예상)},
  title style={font=\small\bfseries}
]
% Theia3D
\addplot[fill=noteorange!60] coordinates {
  (Wind-up, 3) (Stride, 5) (Late Cocking, 10) (Acceleration, 25) (Deceleration, 18) (Follow-through, 8)
};
% ViTPose + EasyMocap
\addplot[fill=cvpurple!40] coordinates {
  (Wind-up, 2.5) (Stride, 4) (Late Cocking, 8) (Acceleration, 20) (Deceleration, 14) (Follow-through, 6)
};
% GART-Pitch
\addplot[fill=defblue!60] coordinates {
  (Wind-up, 1.5) (Stride, 2) (Late Cocking, 3) (Acceleration, 7) (Deceleration, 5) (Follow-through, 3)
};
\legend{Theia3D, ViTPose+EasyMocap, GART-Pitch (Ours)}
\end{axis}
\end{tikzpicture}
\caption{투구 위상별 어깨 관절각 RMSE 예상 결과.
가속 단계에서 기존 방법(25$^\circ$, 20$^\circ$)과 GART-Pitch(7$^\circ$)의 격차가 가장 크다.
이는 광도 기반 자세 정제가 가장 빠른 동작 구간에서 가장 큰 효과를 발휘함을 보여준다.}
\label{fig:phase_accuracy}
\end{figure}


\section{실험 3: 임상 지표 일치도}

\begin{figure}[H]
\centering
\begin{tikzpicture}
\begin{axis}[
  width=10cm, height=6cm,
  xlabel={평균 (Vicon + GART-Pitch) / 2 ($^\circ$)},
  ylabel={차이 (GART-Pitch $-$ Vicon) ($^\circ$)},
  xmin=140, xmax=190,
  ymin=-10, ymax=10,
  grid=both,
  grid style={gray!20},
  title={\small Bland--Altman: 어깨 MER (예상)},
  title style={font=\small\bfseries}
]
% Simulated GART-Pitch data (tighter than existing)
\addplot[only marks, mark=*, mark size=1.5pt, defblue!70] coordinates {
  (150,1.2)(152,-0.8)(155,2.1)(158,-1.5)(160,0.3)(162,1.8)(165,-0.6)
  (167,2.5)(168,-1.2)(170,0.7)(172,1.4)(175,-0.9)(177,0.2)(180,1.6)
  (148,0.5)(153,-1.8)(156,1.0)(159,2.3)(161,-0.4)(163,0.9)(166,-2.0)
  (169,1.5)(171,-0.7)(173,2.0)(176,-1.3)(178,0.8)(182,-0.2)(185,1.1)
};
% Bias
\addplot[red, thick, domain=140:190] {0.4};
\node[red, font=\small] at (axis cs:187,1.2) {bias = 0.4$^\circ$};
% LoA
\addplot[red, dashed, domain=140:190] {0.4+1.96*1.3};
\addplot[red, dashed, domain=140:190] {0.4-1.96*1.3};
\node[red, font=\tiny] at (axis cs:187,3.5) {+1.96 SD};
\node[red, font=\tiny] at (axis cs:187,-2.7) {$-$1.96 SD};
\end{axis}
\end{tikzpicture}
\caption{예상 Bland--Altman 플롯: 어깨 MER.
바이어스 0.4$^\circ$, 95\% LoA $\pm$2.5$^\circ$로
임상적으로 우수한 일치도를 기대한다.}
\label{fig:bland_altman_mer}
\end{figure}

\begin{resultbox}[title={예상 결과: 임상 지표 일치도}]
\begin{center}
\renewcommand{\arraystretch}{1.3}
\begin{tabular}{lcccc}
\toprule
\textbf{지표} & \textbf{RMSE} & \textbf{ICC} & \textbf{Bland--Altman Bias} & \textbf{목표 달성} \\
\midrule
MER & 3--5$^\circ$ & $> 0.95$ & $< 1^\circ$ & 우수 \\
팔꿈치 굴곡 & 2--3$^\circ$ & $> 0.97$ & $< 0.5^\circ$ & 우수 \\
어깨 내회전 속도 & 200--400\degps & $> 0.90$ & $< 100$\degps & 양호 \\
외반 토크 (nRMSE) & 5--10\% & $> 0.85$ & --- & 수용 가능 \\
\bottomrule
\end{tabular}
\end{center}
\end{resultbox}


\section{실험 4: Ablation Study}

\begin{resultbox}[title={예상 결과: Ablation Study (MPJPE, mm)}]
\begin{center}
\renewcommand{\arraystretch}{1.3}
\begin{tabular}{lcc}
\toprule
\textbf{변형} & \textbf{전체 MPJPE} & \textbf{가속단계 MPJPE} \\
\midrule
SMPL-X Only (초기화) & 40--50 & 70--90 \\
+ 3DGS Pose Refinement (상수 $\lambda$) & 25--30 & 40--50 \\
+ Phase-Adaptive $\lambda_{\text{temp}}$ & 18--22 & 22--28 \\
+ Biomech. Constraint & 17--20 & 20--25 \\
\textbf{+ Coarse-to-Fine (Full)} & \textbf{15--20} & \textbf{18--22} \\
\bottomrule
\end{tabular}
\end{center}
\end{resultbox}

\begin{keyidea}
Ablation 결과의 핵심 시사점:
\begin{enumerate}[nosep]
  \item 3DGS 자세 정제만으로도 SMPL-X 대비 $\sim$40\% 개선 ($\rightarrow$ 렌더링 기반 정제의 효과 입증)
  \item 위상별 적응형 정규화가 가속 단계에서 가장 큰 개선 ($\rightarrow$ 투구 특화 설계의 효과)
  \item 생체역학 제약은 전체 정확도보다 해부학적 타당성 향상에 기여
  \item Coarse-to-Fine 전략이 안정적 수렴에 기여
\end{enumerate}
\end{keyidea}


% ====================================================================
% CHAPTER 6: 기술적 도전과 해결 전략
% ====================================================================
\chapter{기술적 도전과 해결 전략}
\label{ch:challenges}

\section{고속 동작에서의 3DGS 한계}

\begin{table}[H]
\centering
\caption{기술적 도전과 대응 전략}
\label{tab:challenges}
\renewcommand{\arraystretch}{1.4}
\begin{tabularx}{\textwidth}{p{3.5cm}X}
\toprule
\textbf{도전} & \textbf{해결 전략} \\
\midrule
가속 30\,ms에서 7프레임 확보 &
인접 프레임 warm-start; 위상별 적응형 시간 정규화;
필요 시 Vicon 고주파 데이터로 보간하여 중간 프레임 의사 GT 생성 \\
\addlinespace
모션 블러 &
1/1000초 이하 노출 시간; 충분한 조명 (1000+ lux);
블러 인식 Gaussian (모션 블러를 모델링하는 확장 가능) \\
\addlinespace
자기 가려짐 (self-occlusion) &
와인드업/팔로스루에서 몸통이 팔을 가림.
7시점의 다양한 각도가 최소 3대 이상 가시 보장.
GART의 정규 공간에서 가려지지 않은 영역이 학습 유지 \\
\addlinespace
계산 비용 &
GART 구조 활용: 캐노니컬 모델 1회 학습 후 프레임별 자세만 최적화.
자세 정제는 자세 파라미터(72개 값)만 갱신하므로 프레임당 $\sim$1--2초.
전체 투구(1--2초, 240--480 프레임): $\sim$15분 내 완료 가능 \\
\addlinespace
손가락/공 세밀 추적 &
SMPL-X의 손 관절(30개) 활용; 공은 별도 트래킹;
릴리스 타이밍 정밀 검출에는 공 위치 변화율 활용 \\
\bottomrule
\end{tabularx}
\end{table}


\section{잠재적 실패 모드와 대비}

\begin{enumerate}[leftmargin=*]
  \item \textbf{3DGS 정제가 국소해에 수렴}: 초기 SMPL-X 피팅이 크게 잘못된 경우,
        광도 손실이 올바른 자세로 안내하지 못할 수 있음.
        $\rightarrow$ Coarse-to-Fine 전략 + 다중 초기화 + 7시점 다시점 제약으로 완화.

  \item \textbf{유니폼의 외형 모호성}: 느슨한 유니폼이 체형을 가려 광도 일치가 불충분할 수 있음.
        $\rightarrow$ 실루엣 손실 + 키포인트 재투영 손실을 보조적으로 활용.

  \item \textbf{시간적 일관성 vs 정확도 트레이드오프}: 과도한 시간 정규화는 빠른 동작을 억제하고,
        부족하면 떨림이 발생.
        $\rightarrow$ 위상별 적응형 가중치가 이 트레이드오프를 관리.
\end{enumerate}


% ====================================================================
% CHAPTER 7: 한계점 및 향후 연구
% ====================================================================
\chapter{한계점 및 향후 연구}
\label{ch:limitations}

\section{현재 연구의 한계}

\begin{enumerate}[leftmargin=*]
  \item \textbf{통제된 환경}: 실내 연구실에서의 촬영만 다루며, 실제 경기장 환경은 고려하지 않음.
        경기장에서는 조명 변화, 관중, 먼 거리 등 추가적 도전이 존재.

  \item \textbf{실시간성 부재}: 프레임당 1--2초의 처리 시간으로 실시간 피드백은 불가능.
        경기 후 분석(post-hoc) 용도에 적합.

  \item \textbf{Vicon 한계의 전이}: Vicon을 gold standard로 사용하지만,
        Vicon 자체도 연조직 아티팩트(STA)로 인한 오차가 존재.
        특히 고속 회전에서 STA가 크며, 이로 인해 ``true'' gold standard와의 비교는 근본적으로 제한됨.

  \item \textbf{피험자 수}: 15--20명은 통계적 검정력에 제한적일 수 있으며,
        다양한 투구 유형(변화구, 사이드암 등)을 포함하지 못할 수 있음.

  \item \textbf{지면 반력 부재}: 역동역학(토크 계산)에 필요한 지면 반력 데이터를
        마커리스 파이프라인에서는 직접 획득할 수 없으며,
        별도의 포스 플레이트가 필요함.
\end{enumerate}


\section{향후 연구 방향}

\begin{enumerate}[leftmargin=*]
  \item \textbf{실시간 GART}: 최근 실시간 3DGS 렌더링 기술(Radl 등, 2024)을
        활용하여 이닝 간 실시간 분석을 목표.

  \item \textbf{소수 카메라 확장}: 7대에서 3--4대로 줄여도 유사한 정확도 유지 가능한지 검증.
        경기장 설치 비용 절감에 직결.

  \item \textbf{이벤트 카메라 결합}: 고속 동작에서 이벤트 카메라의 $\mu$s 시간 해상도를
        활용한 EventSplat 접근법으로 모션 블러 근본 해결.

  \item \textbf{부상 예측 모델}: GART-Pitch로 추출한 정밀 지표를 입력으로
        UCL 부상 예측 모델 구축. Bright 등(2026)의 AUC 0.811을 개선 목표.

  \item \textbf{타 스포츠 확장}: 테니스 서브, 배구 스파이크, 크리켓 볼링 등
        유사한 고속 오버헤드 동작에 적용.
\end{enumerate}


% ====================================================================
% CHAPTER 8: 결론
% ====================================================================
\chapter{결론}
\label{ch:conclusion}

본 연구 제안서는 3D Gaussian Splatting(GART) 기반 마커리스 투구 동작 분석 프레임워크인
\textbf{GART-Pitch}를 제시하였다.

\begin{summary}[title={핵심 논증 요약}]
\begin{enumerate}[leftmargin=*]
  \item \textbf{문제}: 기존 마커리스 SOTA(MPJPE 52--57\,mm)는 투구 분석의 임상 요구를 충족하지 못함.
        특히 어깨 회전 변수(MER, 내회전 속도)에서 가장 큰 오차가 발생하며,
        이는 정확히 투구에서 가장 중요한 변수임.

  \item \textbf{원인}: 기존 방법의 병목은 \textbf{감독 신호의 빈곤}---
        17--33개 키포인트 좌표로 7,000\degps 회전을 정확히 복원하는 것은 본질적으로 부정확.

  \item \textbf{해결}: GART-Pitch는 프레임당 14.5M 픽셀의 광도 일치를 감독 신호로 활용하여
        (63,000$\times$ 더 풍부한 신호), 자세 정확도를 비약적으로 향상.
        투구 위상별 적응형 시간 정규화가 7,000\degps 가속 구간에서도 정확도를 유지.

  \item \textbf{기대 결과}: MPJPE $< 20$\,mm (기존 SOTA 대비 60\%+ 개선),
        MER RMSE $< 5^\circ$, 외반 토크 nRMSE $< 10\%$로
        마커 기반에 근접하는 임상적 정확도 달성.

  \item \textbf{영향}: 마커 없이 Vicon에 근접하는 정확도는
        경기 중 실시간 분석, 과거 영상 소급 분석, 비용 절감, 자연스러운 동작 보장을 가능하게 하여
        투구 생체역학 분석의 패러다임을 전환할 잠재력이 있음.
\end{enumerate}
\end{summary}


% ====================================================================
% APPENDIX
% ====================================================================
\appendix

\chapter{수학적 세부 사항}

\section{3DGS 렌더링 방정식}

깊이 순서로 정렬된 $N$ 개의 Gaussian에 대해, 픽셀 $\vx$의 색상:
\begin{equation}
  C(\vx) = \sum_{k=1}^{N} c_k \, \alpha_k \, G_k(\vx) \prod_{j=1}^{k-1}(1 - \alpha_j \, G_j(\vx))
\end{equation}
여기서 $G_k(\vx)$는 Gaussian $k$의 2D 투영에서 $\vx$ 위치의 값:
\begin{equation}
  G_k(\vx) = \exp\!\Big(-\tfrac{1}{2}(\vx - \vmu_k^{2D})^\top (\mSigma_k^{2D})^{-1}(\vx - \vmu_k^{2D})\Big)
\end{equation}

\section{LBS 변환의 미분}

자세 파라미터 $\vtheta$에 대한 Gaussian 위치의 야코비안:
\begin{equation}
  \frac{\partial \vmu_k^{\text{posed}}}{\partial \vtheta} = \sum_{b=1}^{B} w_{kb} \frac{\partial \mT_b(\vtheta)}{\partial \vtheta} \bar{\vmu}_k^{\text{can}}
\end{equation}
이 미분은 PyTorch의 자동 미분으로 효율적으로 계산된다.


\chapter{용어 사전}

\begin{tabularx}{\textwidth}{lX}
\toprule
\textbf{용어} & \textbf{설명} \\
\midrule
3DGS & 3D Gaussian Splatting. 명시적 3D Gaussian 집합 기반 장면 표현 \\
GART & Gaussian Articulated Representation Templates. SMPL 스키닝을 활용한 동적 인체 3DGS \\
SMPL-X & 인체 파라메트릭 모델 (얼굴, 손 포함) \\
LBS & Linear Blend Skinning. 스켈레톤 기반 메시 변형 \\
MER & Maximum External Rotation. 어깨 최대 외회전 \\
MPJPE & Mean Per Joint Position Error. 3D 관절 위치 오차 (mm) \\
PA-MPJPE & Procrustes Aligned MPJPE. 최적 강체 정합 후 관절 위치 오차 \\
ICC & Intraclass Correlation Coefficient. 측정 일치도 지표 \\
ISB & International Society of Biomechanics. 관절 각도 분해 표준 \\
UCL & Ulnar Collateral Ligament. 팔꿈치 내측 측부 인대 \\
\bottomrule
\end{tabularx}


\end{document}
